\section{Method}
\label{sec:method}

Start with a short overview of \method and point to a method figure, then describe each component in its own subsection.

\subsection{Problem Formulation}
\label{sec:problem}

Define the notation before you use it. For example, a policy $\pi_\theta$ maps an observation $o_t$ and a language instruction $\ell$ to an action chunk $\mathbf{a}_{t:t+H}$:
\begin{equation}
    \mathbf{a}_{t:t+H} \sim \pi_\theta(\cdot \mid o_t, \ell).
    \label{eq:policy}
\end{equation}
Number only the equations you refer to later, and refer to them with \verb|\cref|, e.g.\ \cref{eq:policy}.

\subsection{Cross-references}

Label every section, figure, table, and equation (\verb|\label{sec:...}|, \verb|fig:|, \verb|tab:|, \verb|eq:|) and refer to them with \verb|\cref|, which adds the right name automatically: \cref{sec:intro}, \cref{fig:teaser}, \cref{tab:main}. Use \verb|\Cref| at the start of a sentence.
